\section{Transformer Representation Details}
\label{app:transformer-representations}

Table~\ref{tab:transformer-tokenization-configs} records the concrete
dataset-specific implementations of the three tokenization families used in
Section~\ref{sec:transformer-models}. Here, $d_f$ is the token content
width and $d_p$ is the coordinate width. A token count preceded by
$\leq$ denotes a maximum sequence length; shorter events are masked and
padded only for batching.

\begin{table}[!htbp]
\centering
\caption{Dataset-specific Transformer tokenization configurations. Internal
feature padding, such as the channel padding used by EXO-200 raw patches,
does not introduce additional Transformer tokens.}
\label{tab:transformer-tokenization-configs}
\footnotesize
\setlength{\tabcolsep}{3pt}
\renewcommand{\arraystretch}{1.08}
\begin{tabularx}{\linewidth}{@{}llccX@{}}
\toprule
Dataset
& Family
& Tokens
& $d_f/d_p$
& Construction \\
\midrule

NEXT
& Entity
& $\leq512$
& $2/3$
& Individual detector hits selected under a fixed token cap. Token content
describes the local hit measurement, while centered three-dimensional
coordinates provide position. \\

NEXT
& Region
& $\leq512$
& $2/3$
& Hits are aggregated within occupied three-dimensional voxels of width
$15$ detector-coordinate units. If necessary, the most occupied voxels are
retained under the token cap. \\

NEXT
& Summary
& $\leq512$
& $4/3$
& Hits are ordered using a three-dimensional Morton ordering and partitioned
into balanced spatial groups. Each token summarizes the measurements within
one group. \\

\midrule

MJD
& Entity
& $500$
& $2/3$
& The length-$3800$ waveform is represented by a deterministic mixture
of uniformly spaced and high-importance samples. Local context spans nine
waveform samples. \\

MJD
& Region
& $500$
& $20/3$
& The waveform is divided into $512$ non-overlapping patches containing
$20$ raw amplitudes each. \\

MJD
& Summary
& $500$
& $2/3$
& The same $512$ temporal regions are represented by their segment mean and
root-mean-square amplitude. \\

\midrule

EXO-200
& Entity
& $504$
& $2/6$
& Entities are selected independently within each of six sensor blocks using
equal proportions of uniform and waveform-salient samples. Each token stores
amplitude and nine-sample local RMS. \\

EXO-200
& Region
& $504$
& $150/6$
& Six sensor blocks are divided into seven channel regions and twelve time
regions. Each token retains up to $6\times25$ raw waveform values; smaller
channel regions are padded only within the token feature. \\

EXO-200
& Summary
& $504$
& $4/6$
& The same sensor--temporal regions are described by mean, RMS, maximum
absolute amplitude, and the difference between the second- and first-half
means. \\

\midrule

SuperNEMO
& Entity
& $\leq512$
& $4/3$
& Individual tracker hits form tokens under a deterministic event-level cap,
with centered three-dimensional detector coordinates. \\

SuperNEMO
& Region
& $\leq512$
& $8/3$
& Hits are aggregated into occupied cubic spatial regions with width
$88\,\mathrm{mm}$. Occupancy-based selection is used if an event exceeds the
token cap. \\

SuperNEMO
& Summary
& $\leq16$
& $8/3$
& Hits are Morton ordered and divided into at most sixteen balanced spatial
groups, each represented by aggregated event statistics. \\

\bottomrule
\end{tabularx}
\end{table}

NEXT and SuperNEMO use three-dimensional detector position coordinates. MJD
uses normalized time together with normalized waveform level and local
temporal change. EXO-200 uses normalized time, within-block channel position,
detector side, and indicators for the U-wire, V-wire, and APD sensor
families.

\FloatBarrier